% --------------------------------------------------------------------
%  Singular spectrum of the learning matrix X for an open-boundary TFIM.
%  Requires: tfim_singular_spectrum.csv  (cols: j, sigma)
% --------------------------------------------------------------------
\documentclass[border=4pt]{standalone}

\usepackage{pgfplots}
\usepackage{pgfplotstable}
\usepackage{amsmath,amssymb}
\pgfplotsset{compat=1.18}

\definecolor{royalblue}{HTML}{1f6fb8}
\definecolor{softred}{HTML}{d65c5c}

\begin{document}

\pgfplotstableread[col sep=comma]{tfim_singular_spectrum.csv}\spectrum

\begin{tikzpicture}
\begin{axis}[
    width=11cm,
    height=7cm,
    xlabel={singular-value index $j$ (sorted ascending)},
    ylabel={$\sigma_j(X)$},
    ymode=log,
    log basis y=10,
    grid=both,
    grid style={line width=.1pt, draw=gray!25},
    major grid style={line width=.2pt, draw=gray!50},
    xmin=0,
    enlarge x limits=0.05,
    legend style={font=\small, at={(0.98,0.02)}, anchor=south east,
                  draw=none, fill=white, fill opacity=0.85, text opacity=1},
    title={Open-boundary TFIM,\ $n=10$,\ $t=400$,\ $N_S=10^3$},
    title style={font=\small},
]

\addplot+[only marks, mark=*, mark size=1.6pt, color=royalblue]
    table[x=j, y=sigma]{\spectrum};
\addlegendentry{$\sigma_j(X)$}

% Highlight the smallest singular value (index j=1 after ascending sort).
\addplot+[only marks, mark=*, mark size=2.6pt, color=softred]
    table[x=j, y=sigma, restrict expr to domain={\thisrow{j}}{0:1}]{\spectrum};
\addlegendentry{$\sigma_1$ (null direction $\approx h$)}

\node[anchor=west, font=\scriptsize, color=softred]
    at (axis cs:1.6, 1e-13) {$\sigma_1\approx 10^{-15}$ (machine precision)};

\node[anchor=west, font=\scriptsize, color=royalblue]
    at (axis cs:2.5, 1e-3) {gap $\sigma_2/\sigma_1 \sim 10^{11}$};

\end{axis}
\end{tikzpicture}

\end{document}
